\documentclass[a4paper, 11pt]{article}
\usepackage[utf8]{inputenc}
\usepackage[spanish]{babel}
\usepackage{geometry}
\usepackage{hyperref}

\geometry{a4paper, margin=1in}

\title{Documentación Técnica: Proyecto de Análisis Bibliométrico}
\author{Generado por Gemini}
\date{\today}

\begin{document}

\maketitle

\section{Arquitectura General}

El proyecto sigue una arquitectura de scripts modulares. Cada requerimiento funcional principal está encapsulado en su propio archivo Python. El flujo de datos es secuencial: los primeros scripts se encargan de la limpieza y preparación de datos, y los scripts posteriores consumen esos datos para realizar análisis específicos.

\begin{itemize}
    \item \textbf{Entrada de Datos:} Archivos \texttt{.bib} crudos de las bases de datos ACM y ScienceDirect.
    \item \textbf{Procesamiento Principal:} Scripts de Python que utilizan librerías de ciencia de datos y procesamiento de lenguaje natural (Pandas, Scikit-learn, NLTK, Gensim, Sentence-Transformers).
    \item \textbf{Salida de Datos:} Archivos \texttt{.bib} limpios, archivos de texto (\texttt{.txt}), imágenes de visualización (\texttt{.png}, \texttt{.html}) y un reporte consolidado en PDF.
\end{itemize}

\section{Detalles de Implementación por Requerimiento}

\subsection{Requerimiento 1: Automatización de Datos}
\begin{itemize}
    \item \textbf{Scripts:} \texttt{requerimiento1.py}, \texttt{verificar2.py} (orquestados por \texttt{aplication.py}).
    \item \textbf{Implementación Técnica:}
    \begin{enumerate}
        \item \texttt{concat_bib_files}: Lee todos los archivos \texttt{.bib} de las carpetas de descarga y los concatena en \texttt{todos_raw.bib}.
        \item \texttt{remove_duplicates}: Utiliza la librería \texttt{bibtexparser} para cargar el archivo concatenado. Itera sobre cada entrada, normaliza el campo \texttt{title} (a minúsculas y sin espacios extra) y usa un \texttt{set} de Python para registrar los títulos ya vistos. Esto garantiza una comprobación de duplicados eficiente con una complejidad de O(n). Los artículos únicos y duplicados se guardan en archivos separados.
        \item \texttt{verificar_y_filtrar}: Vuelve a cargar los artículos únicos y verifica que los campos \texttt{title} y \texttt{abstract} no estén vacíos, separando los registros en válidos e inválidos.
    \end{enumerate}
\end{itemize}

\subsection{Requerimiento 2: Similitud Textual}
\begin{itemize}
    \item \textbf{Script:} \texttt{requerimiento2.py}.
    \item \textbf{Implementación Técnica:}
    \begin{enumerate}
        \item \textbf{Levenshtein:} Implementa el algoritmo clásico con una matriz de \texttt{numpy} para calcular el costo de edición entre dos textos.
        \item \textbf{Jaccard:} Convierte los textos a \texttt{set} de palabras y calcula la intersección sobre la unión.
        \item \textbf{TF-IDF \& BoW:} Usa \texttt{TfidfVectorizer} y \texttt{CountVectorizer} de \texttt{scikit-learn} para convertir los textos en vectores. La similitud se calcula con \texttt{cosine_similarity}.
        \item \textbf{Sentence-BERT (SBERT):} Carga el modelo pre-entrenado \texttt{all-MiniLM-L6-v2} de la librería \texttt{sentence-transformers}. Este modelo convierte párrafos enteros en "embeddings" semánticos. La similitud se mide con \texttt{util.cos_sim}.
        \item \textbf{Doc2Vec:} Entrena un modelo \texttt{Doc2Vec} de \texttt{gensim} sobre todo el corpus de resúmenes para aprender vectores que representen cada documento. Luego, infiere los vectores para los dos artículos seleccionados y calcula su similitud.
    \end{enumerate}
\end{itemize}

\subsection{Requerimiento 3: Análisis de Palabras Clave}
\begin{itemize}
    \item \textbf{Script:} \texttt{requerimiento3.py}.
    \item \textbf{Implementación Técnica:}
    \begin{enumerate}
        \item \textbf{Frecuencia Predefinida:} Utiliza expresiones regulares (\texttt{re.findall}) para contar la aparición de cada palabra clave de la lista dada.
        \item \textbf{Generación con TF-IDF:} Emplea \texttt{TfidfVectorizer} de \texttt{scikit-learn} para identificar los términos más relevantes. Se configura para ignorar palabras comunes en inglés y buscar términos de 1 o 2 palabras. Los 15 términos con la puntuación TF-IDF más alta se seleccionan.
    \end{enumerate}
\end{itemize}

\subsection{Requerimiento 4: Agrupamiento Jerárquico}
\begin{itemize}
    \item \textbf{Script:} \texttt{requerimiento4.py}.
    \item \textbf{Implementación Técnica:}
    \begin{enumerate}
        \item \textbf{Vectorización:} Usa \texttt{SentenceTransformer} para convertir los resúmenes en embeddings.
        \item \textbf{Clustering:} Utiliza la función \texttt{linkage} de \texttt{scipy.cluster.hierarchy} con los métodos 'ward', 'complete' y 'average'.
        \item \textbf{Visualización:} La función \texttt{dendrogram} de \texttt{scipy} se usa junto con \texttt{matplotlib} para graficar y guardar los árboles jerárquicos como archivos \texttt{.png}.
    \end{enumerate}
\end{itemize}

\subsection{Requerimiento 5: Análisis Visual}
\begin{itemize}
    \item \textbf{Script:} \texttt{requerimiento5.py}.
    \item \textbf{Implementación Técnica:}
    \begin{enumerate}
        \item \textbf{Carga de Datos:} Usa \texttt{bibtexparser} y \texttt{pandas} para cargar los datos en un DataFrame.
        \item \textbf{Nube de Palabras:} Usa la librería \texttt{wordcloud}.
        \item \textbf{Gráficos Temporales:} Usa \texttt{pandas} para agrupar los datos y \texttt{matplotlib} para generar los gráficos de barras.
        \item \textbf{Mapa de Calor:} Extrae el país del autor, lo normaliza con \texttt{pycountry_convert} y usa \texttt{plotly.express.choropleth} para generar un mapa interactivo \texttt{.html}.
        \item \textbf{Reporte PDF:} Utiliza la librería \texttt{fpdf2} para consolidar las imágenes \texttt{.png} generadas en un único documento.
    \end{enumerate}
\end{itemize}

\end{document}
